\section{Introduction}
\label{sec:introduction}

\subsection{Motivation: what does a constrained agent actually buy?}

LLM-based coding agents have moved from one-shot autocomplete to multi-step autonomous workflows: spec-first design, test-driven implementation, multi-role code review, persistent decision traces. Each procedural constraint exchanges a degree of model freedom for a presumed gain in output predictability or downstream correctness. The exchange is not free: every additional verifier call, every spec-then-test cycle, every persisted artefact adds tokens, latency, and engineering complexity. The empirical question is sharp:
\begin{quote}
\emph{Given a coding task, a model, and an inference budget, does adding procedural structure measurably improve the distribution of outcomes, and at what cost?}
\end{quote}

This paper proposes a measurement methodology---not a general theory---that answers this question on operationally well-defined components: target-correctness, robustness, conformance, and token cost. We deliberately separate \emph{what we can prove formally}, \emph{what we can measure}, \emph{what we conjecture}, and \emph{what we can empirically falsify}. Every numerical claim is reproducible from artefacts in the companion repository.

\subsection{Contributions}

We make four contributions, three positive and one negative.

\begin{enumerate}[leftmargin=2em]
\item A \textbf{minimal formal model} (\Cref{sec:formal-model}) of a procedurally-constrained LLM coding agent at the level of \emph{load-bearing primitives only}: stochastic transducer, verifier kernels (stochastic, accommodating LLM-based reviewers), held-out information constraint $Y_S \perp h_t$ (\Cref{eq:holdout-condition}), behavior signature $B = \varphi(\tau)$, conformance/semantics predicates $\PredC, \PredS$ (\Cref{eq:predC,eq:predS}), and the decisional cost-aware metric $\Vmetric$ with explicit notational disambiguation from $V$. The full CMDP extension (state augmentation, regime modes, escalation map, mode-conditioned action sets) is published as a standalone technical report (\Cref{sec:cmdp-pointer}; \texttt{docs/paper/v1.4/cmdp-extension-techreport.md}) for future work requiring policy-space dynamics; it is \emph{not used} by any empirical result of this paper.

\item A \textbf{decisional metric} $\Vmetric$ (\Cref{sec:decisional-metric}) and an \textbf{eFOSD} validation protocol (\Cref{def:efosd}) that admit a closed-form break-even cost weight $\lambda_{\mathrm{breakeven}}$ below which the protocol pays. Validated on $n=120$ trajectories across $2$ tasks $\times\,4$ frontier models with explicit Wilson confidence intervals and parse-failure sensitivity analyses.

\item Two \textbf{empirical claims with explicit domain restrictions}:
  \begin{itemize}
    \item \textbf{L1 (Asymmetric-Quorum Joint-Miss Bound, finance-restricted)}: validated in $1{,}400$ LLM reviews ($n = 200$ injected bugs), error decorrelation $\bar{\rho}\!: 0.62 \to 0.29$, joint-miss $22.91\% \to 12.28\%$, paired McNemar $p = 0.0013$. \emph{Falsified across domains} in Phase 5 (\texttt{bankcheck.py}, $p \ge 0.5$ in all $5$ contrasts); the cold contract-first reviewer collapses to $11\%$ when contracts are misaligned with the domain.
    \item \textbf{C1 (Cold reviewer transfer)}: contract-aligned specificity replicates across $3$ reviewer families on a $3$-library third-party sample (\Cref{sec:phase10}); library-class is a candidate but unconfirmed moderator (provenance-label experiment in \Cref{sec:phase11} tests this directly). A naturalistic replication on CPython \texttt{csv.py} bugs (\Cref{sec:phase12}, $n = 8$ after strict inclusion) is underpowered at the harvested sample size: $0/3$ families reach the Bonferroni-corrected primary threshold and the Codex result is contaminated by $100\%$ training-set leakage of CPython issue numbers. Naturalistic C1 transfer at low $n$ with the leakage caveat is \emph{reported but not claimed}; a leakage-controlled benchmark is a Phase-13 desideratum.
  \end{itemize}

\item A \textbf{negative result with theoretical and empirical components}: \textbf{C-T3 (Learning-Loop Convergence)} \emph{converges almost surely} in simulation under conservative trigger, bounded injection, and weak suppression, but the LLM-grounded counterpart \emph{shows no improvement} under two pre-registered injection formats (abstract-name and concrete-example, both at $p \in [0.95, 0.97]$ against $H_1$ ``injected $>$ baseline''). This documents a \emph{simulation-to-LLM gap}: convergence properties of pattern-extraction loops do not transfer to measurable detection improvements at realistic $n$.
\end{enumerate}

A consolidated \textbf{ten-mode failure taxonomy} (\Cref{sec:failure-modes}) classifies the regimes under which the framework breaks: uniform-failure collapse (the regime that broke our pilot's Shannon-entropy formulation), oracle leakage, symmetric-reviewer anchoring, model-strength ceiling (Claude-vs-Codex confounder), conformance-without-semantics, off-by-one blind spots, out-of-domain prompt collapse, simulation-to-LLM gap, sample-dependent reviewer-choice sensitivity, and descriptive injection-anchoring effects.

\subsection{Scope and what we do not claim}

We are explicit about scope:
\begin{itemize}[leftmargin=2em]
\item We do not claim that constrained protocols are universally beneficial; the break-even $\lambda_{\mathrm{breakeven}}$ is task- and model-specific.
\item We do not claim L1 generalizes outside the finance contract domain; Phase 5 shows transfer requires re-tooled contracts.
\item We do not claim C1 establishes a general property of cold-reviewer prompts; the absolute-detection rate varies across reviewer families and library classes in ways that are not yet causally attributed.
\item We do not propose new theorems in stochastic dominance, mutation testing, or CMDP theory; we apply established tools as measurement instruments.
\end{itemize}

\subsection{Reproducibility}

Every numerical claim in this paper is recomputed from raw artefacts by a validation script in \texttt{docs/paper/validations/} (Phases 4--6) or in each phase's experiment directory (Phases 7--11). Raw reviewer outputs, AST mutator code, pre-registration documents, and Codex peer-review reports are committed to the repository at the indicated paths. Reviewer prompts, contracts (auto-extracted and mismatched-bankcheck), and reference modules are bit-identical to those used at runtime.

\subsection{Paper structure}

\Cref{sec:related-work} surveys related work on LLM agents, mutation testing, and stochastic dominance applied to code-generation systems. \Cref{sec:formal-model} fixes the load-bearing primitives (stochastic transducer, verifier kernels, held-out condition, behavior signature); the full CMDP extension is deferred to a standalone technical report (\Cref{sec:cmdp-pointer}). \Cref{sec:decisional-metric} introduces $\Vmetric$, defines empirical FOSD, and states Empirical Claim E1. \Cref{sec:asymmetric-verification} states the joint-miss bound (Lemma L1) and reports its empirical mechanism within the finance domain. \Cref{sec:learning-loop} formalizes C-T3 and reports the simulation-to-LLM gap. \Cref{sec:experiments} reports the experimental phases 1--11 with pre-registered metrics and computational validation. \Cref{sec:discussion} catalogues the ten failure modes and discusses external validity, including the cross-phase model-version confounder. \Cref{sec:conclusion} closes with explicit Phase-12 desiderata.
